\begin{longrotatetable}

\begin{longtable}{lrrrrrrrrrrrrl}
\caption{\activityscope\ discoveries, independent recoveries, and candidates. Objects with observed activity are marked with an asterisk.}
\label{tab:activityscope_results}\\
\toprule
Object & $q$ {\footnotesize (au)} & $a$ {\footnotesize (au)} & $e$ & $i~(^\circ)$ &
$T_J$\tabnote{Tisserand parameter with respect to Jupiter.} &
$H_V$\tabnote{Absolute magnitude in the $V$ band, assuming an inert object.} &
$N_{\rm opp}$\tabnote{Number of oppositions on which the object has been observed.} &
{\footnotesize $E[N_{\rm opp}]$}\tabnote{Expected number of observed oppositions for an inert object of the same orbit and $H_V$, predicted by the regression model.} &
{\footnotesize $P(N_{\rm opp}\ge4)$}\tabnote{Binary-classifier probability that an inert object of the same orbit and $H_V$ would be observed on at least four oppositions.} &
$\Delta Q$\tabnote{Quantile deficit $\Delta Q = Q_{0.006} - N_{\rm opp}$, where $Q_{0.006}$ is the model's $0.006$-quantile prediction of the opposition count. Larger positive values indicate a larger deficit.} &
$S_{\rm EV}$\tabnote{Poisson lower-tail anomaly score $S_{\rm EV}=F_{\rm Pois}(N_{\rm opp}-1;\,E[N_{\rm opp}]-1)$. Smaller values indicate a larger deficit relative to $E[N_{\rm opp}]$.} &
$\Delta H$\tabnote{\revision{Magnitude offset required to explain anomaly such that with this offset the dimmer $H_V$ would yield $E[N_{\rm opp}]$ within 0.5 of actual $N_{\rm opp}$.}}&
Note \\
\midrule
\endfirsthead

\caption[]{\activityscope\ discoveries, independent recoveries, and candidates. Objects with observed activity are marked with an asterisk. \mbox{(continued)}}\\
\toprule
Object & $q$ {\footnotesize (au)} & $a$ {\footnotesize (au)} & $e$ & $i~(^\circ)$ &
$T_J$\textsuperscript{a} & $H_V$\textsuperscript{b} & $N_{\rm opp}$\textsuperscript{c} &
{\footnotesize $E[N_{\rm opp}]$}\textsuperscript{d} & {\footnotesize $P(N_{\rm opp}\ge4)$}\textsuperscript{e} &
$\Delta Q$\textsuperscript{f} & $S_{\rm EV}$\textsuperscript{g} & $\Delta H$\textsuperscript{h} & Note \\
\midrule
\endhead

\midrule
\multicolumn{14}{r}{Continued on next page}\\
\endfoot

\endlastfoot

\multicolumn{14}{l}{\textbf{Cometary orbits (multiple oppositions when flagged)} \textnormal{\footnotesize (sorted by $S_{\rm EV}$, ascending)}}\\
\midrule

\objectcell{2010~RH$_{69}$$^{\ast}$$^{\ddag}$} &
3.90 & 4.27 & 0.087 & 11.92 & 2.97 & 13.9 & 3 & 18.4 & 0.999997 & 11.8 & 0.00000 & 2.4 &
\notecell{Comet independently re-flagged by this work} \\

\objectcell{362P/(457175) 2008~GO$_{98}$$^{\ast}$} &
2.86 & 3.97 & 0.279 & 15.56 & 2.93 & 12.9 & 7 & 24.1 & 0.999996 & 9.6 & 0.00003 & 3.2 &
\notecell{\citet{2017CBET.4411....1L}} \\

\objectcell{282P/(323137) 2003~BM$_{80}$$^{\ast}$} &
3.44 & 4.23 & 0.188 & 5.81 & 2.99 & 13.6 & 7 & 22.2 & 0.999998 & 9.6 & 0.00011 & 2.1 &
\notecell{\citet{2022ApJ...937L...2C,chandlerActiveAsteroidsCitizen2024}} \\

\objectcell{2017~QN$_{84}$$^{\ast}$} &
2.47 & 3.77 & 0.343 & 12.07 & 2.94 & 15.0 & 4 & 12.8 & 0.999998 & 4.9 & 0.00261 & 2.0 &
\notecell{\citet{chandlerActiveAsteroidsCitizen2024,2025epsc.conf.1661F}} \\

\objectcell{2008~VK$_{110}$$^{\Vert}$} &
4.44 & 7.17 & 0.381 & 7.39 & 2.88 & 13.8$^{\S}$ & 3 & 7.6 & 0.986391 & -0.2 & 0.04190 & 1.2 &
\notecell{Identified by \activityscope{}; Possible morphologic confirmation; Brightness varies across oppositions} \\

\midrule
\multicolumn{14}{l}{\textbf{Cometary orbits (single opposition when flagged)} \textnormal{\footnotesize (sorted by $P(N_{\rm opp}\ge4)$, descending)}}\\
\midrule

\objectcell{P/2025~UX$_{109}$ (Ye)$^{\ast}$} &
2.57 & 3.81 & 0.325 & 3.18 & 2.98 & 15.0 & 1 & 16.7 & 0.999999 & 11.8 & 0.00000 & 4.5 &
\notecell{Independent identification by \activityscope\tabnote{P. VanWylen, \url{https://groups.io/g/comets-ml/message/34359}.}} \\

\objectcell{2001~BV$_{70}$} &
2.06 & 3.59 & 0.426 & 4.33 & 2.95 & 15.1 & 1 & 16.0 & 0.999999 & 10.9 & 0.00000 & 4.4 &
\notecell{High-confidence candidate; archival confirmation unavailable} \\

\objectcell{2019~OE$_{31}$$^{\ast}$} &
3.94 & 4.38 & 0.100 & 5.22 & 3.01 & 14.6 & 1 & 16.8 & 0.999997 & 11.9 & 0.00000 & 3.5 &
\notecell{\citet{oldroydCometaryActivityDiscovered2023}} \\

\objectcell{2018~BJ$_{11}$$^{\ast}$} &
3.25 & 4.20 & 0.226 & 3.43 & 2.99 & 15.6 & 1 & 10.4 & 0.999994 & 5.0 & 0.00009 & 3.3 &
\notecell{Independent identification; Activity visible in archival imagery\tabnote{This object was on the Possible Comet Confirmation Page prior to our reconfirmation and discovery of archival imagery.}} \\

\objectcell{P/2022~BV$_{9}$ (Lemmon)$^{\ast}$} &
3.35 & 4.37 & 0.233 & 11.92 & 2.93 & 13.8 & 1$^{\dagger}$ & 18.6 & 0.999985 & 14.0 & 0.00000 & 4.1 &
\notecell{Identified by \activityscope\tabnote{P. VanWylen, \url{https://groups.io/g/comets-ml/message/31369}.}} \\

\objectcell{489P/Denning$^{\ast}$} &
1.56 & 4.42 & 0.647 & 4.03 & 2.58 & 15.3 & 1$^{\dagger}$ & 8.5 & 0.999472 & 3.3 & 0.00054 & 3.9 &
\notecell{Recovered after 130 yr; initial identification by \activityscope\tabnote{P. VanWylen, \url{https://groups.io/g/comets-ml/message/32365}.}} \\

\objectcell{2024~XE$_{22}$$^{\ast}$} &
4.54 & 5.87 & 0.227 & 27.16 & 2.73 & 14.3 & 1$^{\dagger}$ & 6.3 & 0.989985 & 1.3 & 0.00483 & 3.4 &
\notecell{Identified by \activityscope{}; activity observed in 2025 November} \\

\objectcell{2009~DP$_{2}$$^{\ast}$} &
3.82 & 6.65 & 0.425 & 26.77 & 2.61 & 14.2$^{\S}$ & 1$^{\dagger}$ & 5.4 & 0.936707 & 1.0 & 0.01267 & 4.2 &
\notecell{Identified by \activityscope{}; recovered active in 2026} \\

\midrule
\newpage
\multicolumn{14}{l}{\textbf{Asteroidal orbits} \textnormal{\footnotesize (sorted by $P(N_{\rm opp}\ge4)$, descending)}}\\
\midrule

\objectcell{P/2023~JN$_{16}$ (Lemmon)$^{\ast}$} &
2.30 & 2.70 & 0.147 & 3.70 & 3.35 & 15.7 & 1$^{\dagger}$ & 19.7 & 0.999999 & 15.0 & 0.00000 & 4.5 &
\notecell{Initial identification by \activityscope\tabnote{P. VanWylen, \url{https://groups.io/g/mpml/message/38758}.}; \citet{2024MPEC....Q....4.,2024CBET.5430....1A,2024DPS....5620804T,2025epsc.conf..471A,2025AcA....75..343B}} \\

\objectcell{2025~HV$_{38}$} &
2.37 & 2.48 & 0.048 & 7.24 & 3.46 & 15.4 & 1 & 21.1 & 0.999999 & 15.5 & 0.00000 & 5.2 &
\notecell{Exceptionally strong candidate; not yet morphologically confirmed; not observed in 2026 despite $V \approx 18.3$} \\

\objectcell{2015~BC$_{566}$$^{\ast}$} &
2.95 & 3.06 & 0.037 & 11.69 & 3.20 & 16.3 & 1 & 13.7 & 0.999999 & 8.7 & 0.00000 & 2.9 &
\notecell{Identified by \activityscope; Active in DECam and MegaCam} \\

\objectcell{2007~VB$_{146}$} &
2.44 & 2.75 & 0.114 & 3.49 & 3.33 & 16.9 & 1 & 14.8 & 0.999999 & 10.0 & 0.00000 & 2.7 &
\notecell{Negative archival recovery; very high confidence candidate} \\

\objectcell{2008~BJ$_{22}$$^{\ast}$} &
2.94 & 3.07 & 0.042 & 11.51 & 3.20 & 15.2 & 1 & 20.6 & 0.999999 & 15.4 & 0.00000 & 4.3 &
\notecell{Initial identification by \activityscope\tabnote{S. Deen, \url{https://groups.io/g/comets-ml/message/31013}.}; \citet{2023RNAAS...7...48D}} \\

\objectcell{2025~VZ$_{8}$$^{\ast}$} &
2.39 & 2.62 & 0.086 & 12.10 & 3.37 & 16.8 & 1 & 12.6 & 0.999998 & 7.2 & 0.00001 & 3.1 &
\notecell{Initial identification by \activityscope\tabnote{P. VanWylen, \url{https://groups.io/g/mpml/message/41136}.}; \citet{2026RNAAS..10....1H}} \\

\objectcell{2021~AY$_{8}$$^{\ast}$} &
2.57 & 2.72 & 0.057 & 5.01 & 3.35 & 17.5 & 1$^{\dagger}$ & 12.3 & 0.999997 & 7.0 & 0.00001 & 3.0 &
\notecell{Identified by \activityscope; Confirmed active in DECam and Pan-STARRS archives} \\

\objectcell{2009~FP$_{8}$} &
2.48 & 3.01 & 0.178 & 4.63 & 3.22 & 16.9 & 1 & 11.8 & 0.999995 & 6.6 & 0.00002 & 3.6 &
\notecell{Negative archival recovery; very high confidence candidate} \\

\objectcell{2010~TR$_{241}$} &
2.75 & 3.05 & 0.100 & 12.42 & 3.19 & 17.0 & 1 & 10.6 & 0.999994 & 5.3 & 0.00007 & 2.0 &
\notecell{Negative archival recovery; high confidence candidate} \\

\objectcell{P/2017~FL$_{36}$ (PANSTARRS)$^{\ast}$} &
2.83 & 2.93 & 0.034 & 15.70 & 3.22 & 16.6 & 1 & 9.5 & 0.999980 & 4.1 & 0.00020 & 2.8 &
\notecell{Initial identification by \activityscope\tabnote{P. VanWylen, \url{https://groups.io/g/comets-ml/message/33447}.}; \citet{2025CBET.5551....1L,2025MPEC....H...86C}} \\

\objectcell{2002~CW$_{116}$$^{\ast}$} &
2.07 & 2.69 & 0.231 & 8.14 & 3.32 & 17.5 & 1 & 9.2 & 0.999881 & 3.9 & 0.00029 & 2.9 &
\notecell{Identified by \activityscope; confirmed active in SDSS archives\tabnote{P. VanWylen, S. Deen \url{https://groups.io/g/mpml/topic/104468465}}} \\

\bottomrule
\multicolumn{14}{l}{\footnotesize Unless otherwise noted, all orbital elements are osculating elements from the epoch provided in \texttt{MPCORB}.}\\
\multicolumn{14}{l}{\footnotesize $^{\ast}$\,Observationally confirmed active object.}\\
\multicolumn{14}{l}{\footnotesize $^{\Vert}$\,Activity supported by combination of observational and photometric evidence.}\\
\multicolumn{14}{l}{\footnotesize $^{\dagger}$\,Single-opposition at the time of flagging; tabulated $N_{\rm opp}$, $\Delta Q$, and $S_{\rm EV}$ are the values the model assigns to an otherwise-identical single-opposition object.}\\
\multicolumn{14}{l}{\footnotesize $^{\ddag}$\,Indicates that provided orbital information is from osculating elements taken from the discovery epoch instead of latest available epoch.}\\
\multicolumn{14}{l}{\footnotesize $^{\S}$\,\revision{A multi-opposition object with significant differences in $H_V$ across oppositions. We use the brightest $H_V$.}}\\
\printtabnotes

\end{longtable}

\end{longrotatetable}